\documentclass[12pt]{article}
\usepackage[english]{babel}
\usepackage[utf8x]{inputenc}
\usepackage[T1]{fontenc}
\usepackage{scribe}
\usepackage{listings}


\begin{document}

Songyu Ye 

sy459@cornell.edu

\today
\section{Universal enveloping algebra and PBW}
Consider the tensor algebra of a vector space $E$ \begin{align*}
    T(E) = \bigoplus_{i=0}^\infty E^{\otimes n} 
\end{align*}

\begin{definition}
The \textbf{symmetric algebra} of $E$ is the associative algebra $S(\mf g)$ defined as the quotient \begin{align*}
    S(E) = T(E) / J
\end{align*} 
\end{definition}
where $J = \langle x\otimes y - y\otimes x\rangle$. There are distinguished subspaces $S^n(E) = E^{\otimes n} / (E^{\otimes n} \cap J)$ which satisfies the following \textbf{universal property}. Every symmetric multilinear map to a vector space $E\times \dots \times E\to U$ factors uniquely through $E\times\dots \times E\to S(E)$.

\hfill

$S(E)$ satisfies the following universal property. Every linear map $E\to A$ where $A$ is \red{\textbf{commutative} associative algebra with identity} factors uniquely through $i:E\to S(E)$.

\begin{definition}
The \textbf{exterior algebra} of $E$ is the associative algebra $\Lambda(E)$ defined as the quotient \begin{align*}
    \Lambda(E) = T(E) / I
\end{align*} where $I$ is the two sided ideal generated by all symbols of the form $x\otimes x$ where $x\in E$. There are distinguished subspaces $\Lambda^n(E) = E^{\otimes n} / (I\cap E^{\otimes n})$ so that \begin{align*}
    \Lambda(E) = \bigoplus_{i=0}^\infty \Lambda^n(E)
\end{align*}
\end{definition} which satisfy the following universal property. Every alternating multilinear map to a vector space $E\times \dots \times E \to U$ factors uniquely through through $E\times\dots \times E\to \Lambda(E)$.

\hfill 

$\Lambda(E)$ satisfies the following universal property. Every linear map $\ell: E\to A$ with $\ell^2(v) = 0$ for all $v\in E$ where $A$ is \red{associative algebra with identity} factors uniquely through $i:E\to S(E)$.

\begin{definition}
The \textbf{universal enveloping algebra} of $\mf g$ is the associative algebra $U(\mf g)$ defined as the quotient \begin{align*}
    U(\mf g) = T(\mf g) / \langle x\otimes y - y\otimes x - [x,y]\rangle
\end{align*} Observe that when $\mf g$ is abelian we have that $U(\mf g)  = S(\mf g)$ the symmetric algebra. We have the distinguished subspace $U^n(\mf g) = \pi(\bigoplus_{i\leq n}g^{\otimes n})$ where $\pi:T(\mf g)\to U(\mf g)$ the quotient map. 

\hfill 

$U(\mf g)$ has the universal property that every linear map $f:\mf g\to A$ with $f([x,y]) = f(x)f(y) - f(y)f(x)$ factors uniquely through $i:\mf g\to U(\mf g)$. Note that we have not yet shown that $i$ is injective. This follows from the following theorem.

\begin{theorem}
\textbf{(Poincare-Birkhoff-Witt):}
Let $x_1,\dots,x_n$ denote an ordered
\end{theorem}
\begin{proof}
\red{fill in the proof for this}
\end{proof}
\begin{example}
    Compute the PBW basis in the cases of $\mf u(2,\C),\mf{sl}_2(\C),\mf u(1,1)$. Compute the Casimir element as well.
\end{example}
\end{definition}

\section{Semisimple Lie Algebras}
\begin{proposition}
    If $\mf g$ is a finite dimensional Lie algebra over $K$, then we have a symmetric invariant bilinear form $k:\mf g\times\mf g\to K$ given by \begin{align*}
        k(x,y) = \tr(\ad x \circ \ad y)
    \end{align*}
\end{proposition}
\red{Is the killing form necessarily nonzero? I guess there is Cartan's criterion which says that $k$ is nondegenerate if and only if $\mf g$ is semisimple, which means by definition that $\mf g$ has the radical (maximal solvable ideal) is $0$. This is (in the finite dimensional case) equivalent to the statemnet that $\mf g$ is a direct sum of simple lie algebras.}
\begin{proposition}
    If moreover we assume $\mf g$ is simple (i.e. has no nontrivial ideals), then $k$ is nondegenerate and unique up to scale.
\end{proposition}
\red{what I wrote above was alternatively written as:}
\begin{proposition}
Let $\mf g$ be a finite-dimensional simple Lie algebra over a field $K$. Then there exists, up to scalar, at most one invariant bilinear form on $\mf g$. If a nonzero invariant bilinear form exists it is nondegenerate and symmetric.
\end{proposition}
Why was it written like this? We know it exists right? It is just the Killing form, unless I am mistaken? We just have to worry about the nonzero-ness of the Killing form.
\begin{proposition}
Suppose that $\mf g$ has a \textbf{nondegenerate} invariant bilinear form $B$. Let $x_1,\dots,x_d$ be a basis of $\mf g$ and let $y_1,\dots,y_d$ be the dual basis with respect to the Killing form, i.e. $B(x_i,y_j) = \delta_{ij}$. Then $\Delta = \sum_{i}x_iy_i\in U(\mf g)$ is in the center of $U(\mf g)$.
\end{proposition}
\red{$\Delta$ is called the Casimir element for $\mf g$ with respect to $B$}
\red{example of when we have a bunch of invariant bilinear forms getting us a bunch of different Casimirs?}

\section{Reps of sl2}
\begin{example}
Consider the adjoint representation of $\so_3$ on $\so_4$. We have \begin{align*}
    \ad_{\so_3}(x\in\so_4) &= \begin{bmm}
        A & 0\\
        0 & 0 
    \end{bmm} \begin{bmm}
        B & w^T\\
        -w^T & 0 
    \end{bmm} -  \begin{bmm}
        B & w^T\\
        -w^T & 0 
    \end{bmm}\begin{bmm}
        A & 0\\
        0 & 0 
    \end{bmm} \\
    &= \begin{bmm}
        [A,B] & Aw^t \\
        wA^t & 0
    \end{bmm} = \so_3 \oplus \C^3
\end{align*}
In fact we claim that $\so_3\cong \C^3$, that is the adjoint representation of $\so_3$ is isomorphic to the standard representation. One can explicitly compute the intertwining element to see this isomorphism.
\end{example}
\begin{example}
Consider the adjoint representation of $\sl_2$ on $\sl_3$. We have \begin{align*}
    \sl_3 = \set{\begin{bmm}
        B & v\\
        w^T & -\tr B
    \end{bmm}} 
\end{align*} where $B\in\gl_2$
\end{example}

\section{Borels and Complexification}
Consider $K = U(n)$ which is a maximal compact subgroup of $G = \GL_n\C$. Moreover $G$
is the complexification of $K$.

\begin{definition}
    Let $K$ Lie Group. A complexification of $K$ is a Lie group homomorphism $K\hookrightarrow G$
    where $G$ is a complex Lie group so that $\mf k$ is a real form of $\mf g$, i.e. $\mf k_\C\cong \mf g$.
\end{definition}

Not every real Lie group has a complexification. The standard example of this 
is the universal cover of $\SL_2$. This motivates the following definition: \begin{definition}
    Let $K$ be a connected Lie group. Then the universal complexification of $K$
    is a Lie group homomorphism $K\to G$ where $G$ is a complex Lie group so that
    any other Lie group homomorphism $K\to H$ where $H$ is a complex Lie group factors
    uniquely through $G\to H$, and moreover this map is analytic.
\end{definition}

\begin{theorem}
    Let $K$ be a connected Lie group. Then the universal complexification of $K$
    exists and is unique up to unique isomorphism.
\end{theorem}

\begin{theorem}
    $\SL_n\R$ has complexification $\SL_n\C$.
\end{theorem}

\begin{proof}
    Consider $f:\SL_n\R\to H$ where $H$ is complex analytic. The differential gives us a map 
    $df:\mf{sl}_n\R\to \mf h$ where $\mf h$ is a complex Lie algebra. By the property of 
    complexification of Lie algebras, we get an induced map $\mf{sl}_n\C\to \mf h$. 
    Since $\mf{sl}_n\C$ is simply connected this map is the differential of a unique map $\SL_n\C\to H$.
    \red{there are two things that I need to check here}

    Finally we need to see that $f$ is analytic. This is true because of the following diagram:
\[\begin{tikzcd}
	\sl_n\C & Lie(H) \\
	\SL_n\C & H
	\arrow[from=1-1, to=1-2]
	\arrow[from=1-2, to=2-2]
	\arrow[from=2-1, to=2-2]
	\arrow[from=1-1, to=2-1]
\end{tikzcd}\] We know that every map in the diagram is holomorphic except for perhaps the bottom one. Then we are done
since $\exp$ is a local homeo and near $g\in G$ we can translate $F = L(f(g))\circ F\circ L(g^{-1})$.\end{proof}

\begin{theorem}
    Let $K$ be a compact connected Lie group. Then $K$ has a complexification $K\to G$. This map induces 
    an isomorphism on $\pi_1$. Moreover $\Lie(G) = \Lie(K)\otimes_\R\C$. Any faithful complex 
    representation of $K$ extends to a complex analytic representation of $G$. Any analytic rep of $G$ is completely reducible.
\end{theorem}

\begin{theorem}
    Let $K$ compact connected. Then the universal complexification of $K$ coincides with 
    the complexification of $K$. In particular compact groups have complexifications.
\end{theorem}


In this following discussion we will be considering the following subgroups.
Consider maximal torus $T\subset K$ consising of diagonal matrices with $\abs{\lambda_i} = 1$. 
We have the complex torus $T_\C$ which factors $T_\C = TA$ where $A$ is those
diagonal matrices with $\lambda_i\in\R_{>0}$. We also have \begin{align*}
    G \supset B \supset B_0 \supset N
\end{align*} where $B$ is upper triangular matrices, $B_0$ is upper triangular matrices 
with $\R^+$ on the diagonal (i.e. the connected component of $B$ containing $I$),
 and $N$ is the upper triangular matrices with $1$ on the diagonal.

\begin{proposition}
    (Iwasawa decomposition)
    Every $g\in G$ factors uniquely $g = avk$ where $a\in A$, $v\in N$, $k\in K$. Moreover the map 
    $A\times N\times K\to G$ is a diffeomorphism.
\end{proposition}

\begin{proof}
    Let $v_1,\dots,v_n$ rows of $g$. Gram Schmidt says there exists $t_{ij}$ for $i<j$ so that \begin{align*}
        v_n,v_n+t_{n-1,n-1}v_{n-1},\dots,v_n + t_{1,n-1}v_{n-1} + \dots + t_{1,1}v_1
    \end{align*} are orthogonal. Let $u_i$ be those vectors.
    Put \begin{align*}
        v^{-1} &= \begin{bmatrix}
            1 & t_{1,1} & \dots & t_{1,n-1}\\
            & 1 & \dots & t_{2,n-1}\\
            & & \ddots & \vdots\\
            & & & 1
        \end{bmatrix}\\
        a &= \begin{bmatrix}
            \abs{u_1} & & 0\\
            & \ddots & \\
            0 & & \abs{u_n}
        \end{bmatrix}
    \end{align*}
    Then $k = a^{-1}v^{-1}g$ has orthonormal rows and hence is unitary. This shows onto. 
    Injectiivty follows from the fact that $B_0\cap K = 1$ and $A\cap N = 1$.
\end{proof}

The following lemma will help us prove Lie's theorem.
\begin{lemma}
    (Dynkin) Let $V$ finite dimeisonal over $F$ characteristic $0$. Let $\mf g\subset \gl(V)$ and $\mf h\leq \mf g$ ideal.
    Let $\lambda:\mf h\to F$ linear form. Then \begin{align*}
        W = \set{v\in V: hv = \lambda(h)v\text{ for all }h\in\mf h}
    \end{align*} is $\mf g-$invariant.
\end{lemma}
\begin{proof}
    Let $v\in W$ nonzero and $X\in\mf g$. We want to show $Xv\in W$.
    
    Consider $W_0 = \ideal{v_0,Xv_0,\dots, X^{d-1}v_0}$. Let $Z\in\mf h$. Then $Z(W_0)\subset W_0$ we have the following trace formula \begin{align*}
        \tr Z\vert_{W_0} = \dim(W_0)\lambda(Z)
    \end{align*} This follows from the fact that the matrix for $Z$ with respect to the given basis for $W_0$ is upper triangular. One can see this by induction.
    \begin{align*}
        ZX^iv = XZX^{i-1}v + [Z,X]X^{i-1}v = \lambda(Z)X^iv + \sum_{j<i}c_jX^jv +\lambda([Z,X])X^{i-1}v
    \end{align*} by induction. Then the trace formula clearly follows. Then we have \begin{align*}
        ZXv = XZv + [Z,X]v = \lambda(Z)Xv + \lambda([Z,X])v  = \lambda(Z)Xv
    \end{align*} since $[Z,X]\in\mf h$, so $0 = \dim(W_0)\lambda([X,Z])$ and we are over characteristic zero.
\end{proof}

\begin{theorem}
    (Lie) Let $\mf b\subset \mf{gl}(V)$ solvable where $V$ finite dimensional over $F$ characteristic $0$. \begin{enumerate}
        \item There exists a simultaneous eigenvector for all of $\mf b$.
        \item There exists a basis for $V$ so that $\mf b$ is upper triangular.
    \end{enumerate}
\end{theorem}
\begin{proof}
    Since $\mf b$ is solvable, $[\mf b,\mf b]$ is proper and hence $\mf b$ has a codimension $1$ ideal $\mf h$.
    Let $\lambda$ the corresponding weight and let $W = \set{v\in V: hv = \lambda(h)v\text{ for all }h\in\mf h}$.
    Then $W$ is nonempty and hence $\mf b-$invariant by the lemma. Let $Z\in \mf b\backslash\mf h$.
    Then $Z$ has an eigenvector $v_1$ on $W$ since $F$ is algebraically closed. $v_1$ is an eigenvector for all of $\mf b$ since it already is for $\mf h$.
    Claim (ii) follows from induction on $V/Fv_1$.
\end{proof}

Now let $K$ be a compact connected Lie group, $G$ its complexification. Let $\mf k = \Lie(K)$ and $\mf g = \Lie(G)$. $K$ has a maximal torus $T$,
its complexification $T_\C$ embeds in $G$ by the universal property.

Let $\Phi$ root system of $K$ and pick system of positive roots $\Phi^+$. We will see that this amounts to picking a Borel subgroup. 
We can put \begin{align*}
    \mf n = \bigoplus_{\alpha\in\Phi^+}\mf g_\alpha
\end{align*} where $\mf g_\alpha$ is the $1$-dimensional root space corresponding to $\alpha$. Then $\mf n$ is nilpotent and hence solvable.

\begin{definition}
    A Lie algebra $\mf n$ is nilpotent if its lower central series defined by \begin{align*}
        \mf n_1 = \mf n,\quad \mf n_{k+1} = [\mf n_k,\mf n]
    \end{align*} eventually terminates. This is equivalent to asking for the existence of a chain of ideals \begin{align*}
        \mf n = \mf n_1\supset \mf n_2\supset \dots \supset \mf n_k = 0
    \end{align*} where $\mf n_i \supset [\mf n_{i-1},\mf n]$.
\end{definition} 

\begin{proposition}
    $\mf n$ is nilpotent.
\end{proposition}
\begin{proof}
    Let $\Phi^+_k$ denote the set of positive roots expressible as the sum of at least $k$ positive roots.
    Then $\mf n_k = \bigoplus_{\alpha\in\Phi^+_k}\mf g_\alpha$ is nilpotent by induction on $k$.
\end{proof}
Let $\mf t = \Lie(T)$ and $\mf b = \mf t\oplus \mf n$. Then $\mf b$ is a complex Lie algebra. 
Moreover $[\mf b,\mf b]\subset \mf n$ so $\mf b$ is solvable.

\begin{theorem}
    Put $N = \exp(\mf n)$ and $B = \exp(\mf b)$. Then \begin{enumerate}
        \item $B$ and $N$ are closed subgroups of $G$, and $Lie(B) = \mf b$ and $Lie(N) = \mf n$.  
        \item We can embed $G\hookrightarrow\GL_n\C$ so that $B$ is the subgroup of upper triangular matrices,
        $K$ is the subgroup of unitary matrices, and $T_\C$ is the subgroup of diagonal matrices.
        \item If $\mf v,\mf w$ are complex Lie subalgebras of $\mf n$ so that $\mf n = \mf v\oplus \mf w$ then $V = \exp \mf v$ is a closed Lie subgroup of $G$
        and $N = V\cdot W$ and $V\cap W = 1$.
    \end{enumerate}
\end{theorem}
From this theorem, one obtains the Iwasawa decomposition for all $K\to G$ complexification of compact connected Lie group.

\section{Compact Lie groups}
\begin{theorem}
    (Cartan/Polar decomposition) The multiplication map $P\times U(n)\to \GL_n$ is a diffeomorphism, where $P$
    is the space of positive definite Hermitian matrices.
\end{theorem} This result follows from discussions about the universal covers of these groups. 
\begin{theorem}
    \begin{enumerate}
        \item (Lie's homomorphism theorem) If $G$ is simply connected then every $\mf g\to \mf h$ extends uniquely to $G\to H$.
        \item There exists a simply connected Lie group $G$ with Lie algebra $\mf g$ for all Lie algebras $\mf g$. 
        \item Every rep of $G$ actually comes from a rep of $\tilde G$ where $\tilde G$ is the universal cover.
    \end{enumerate}
\end{theorem}
Say that you want to study the classification of compact Lie groups. One can do this as follows.
\begin{fact}
    If $G$ and $H$ are compact and simply connected and they have the same Lie algebra, then $G\cong H$.
\end{fact}
\begin{fact}
    If $G$ is any Lie group then its universal cover $\tilde G$ is naturally a Lie group for which 
    $G = \tilde G/Z$ where $Z\subset Z(\tilde G)$.
\end{fact}
These two facts imply that to classify compact Lie groups, it suffices to classify
compact simply connected Lie groups, and then compute the centers. These in turn are in 
correspondence with complex semisimple Lie algebras. If I have a compact simply conneced Lie group, I can take
the complexification of its Lie algebra. Conversely, every complex semisimple Lie algebra has a unique compact real form which
is isomorphic to the Lie algebra of a unique compact simply connected Lie group.
\end{document}